\section{Related Work}

\subsection{Query Expansion and Hybrid Retrieval}

BM25 provides a strong lexical first-stage retriever for matching
term evidence in long narratives \citep{robertson2009probabilistic}.
Query expansion can complement this lexical anchor: Query2Doc uses a
large language model to generate a pseudo-document for expansion
\citep{wang2023query2doc}, while related work studies prompted query
expansion and iterative refinement \citep{jagerman2023queryexpansion,chen2024analyze}.
We use the original narrative and expanded formulations as separate
routes, so an expanded query does not replace the lexical anchor.

Hybrid retrieval also combines lexical and semantic signals at
different stages. Dense Passage Retrieval learns a dense first-stage
representation \citep{karpukhin2020dpr}, late interaction retains
token-level matching for efficient passage search
\citep{khattab2020colbert}, and SPLADE makes lexical expansion part of
the representation \citep{formal2021splade,formal2021spladev2}. Our
candidate pool likewise combines complementary retrieval signals; we
evaluate their route-level combination and subsequent serialization
on the development set.

\subsection{Neural Reranking and Rank Fusion}

Rank fusion is useful when component rankings provide complementary
evidence. Reciprocal Rank Fusion is a simple robust baseline for
combining rankings without requiring a common score scale
\citep{cormack2009rrf}. Our Retrieval system applies weighted rank
fusion at route-combination stages and again when combining the
BM25, cross-encoder, and dense-ranking head signals. A cross-encoder tail
rank is then fused with candidate-pool order only over the specified
tail, leaving the head and the variable-depth cutoff stages separable.
BEIR shows that zero-shot behavior varies substantially
across heterogeneous domains and that some neural rerankers can trail
BM25 on particular out-of-domain collections \citep{thakur2021beir}.
This is not a universal weakness of cross-encoders: scaling studies
find that sufficiently large cross-encoders can generalize strongly
and outperform comparable bi-encoders \citep{rosa2022crossencoders}.
Our protected-head choice is based on development-set observations for
ClimbMix, not on a general claim that cross-encoders degrade out of
domain.

Choosing how deep to read a ranked list is itself an open problem
separate from how the list is ordered. Early work optimizes a
per-query cutoff from the shape of the query's own score distribution
\citep{arampatzis2009stopreading}; more recent work selects a
long-context QA read depth from the steepest drop in a sorted
similarity-score distribution, without additional tuning or iterative
model calls \citep{taguchi2025adaptivek}. Our variable-depth
serialization (\S\ref{sec:retrieval-method}) is a relative-threshold
instance of this same family, applied to the score used immediately
before deep-tail reranking rather than to a single-query similarity
distribution.

\subsection{Evidence-Grounded and Agentic RAG}

Retrieval-augmented generation couples evidence selection with answer
generation \citep{lewis2020rag}. REALM and Atlas demonstrate how
retrieval can be integrated with language-model reasoning and few-shot
adaptation \citep{guu2020realm,izacard2023atlas}; WebGPT similarly
studies browser-assisted question answering with feedback
\citep{nakano2021webgpt}. Our RAG system also acquires evidence before
generation, but it limits additional reading and handles citation
ordering in a separate finalization step.
Self-RAG addresses when to retrieve by training a single model to emit
reflection tokens that decide, and critique, retrieval and generation
end to end \citep{asai2024selfrag}; our sufficiency judge makes the
analogous retrieve-more-or-stop decision but as an explicit, bounded,
inference-time policy rather than a trained-in behavior. The resulting
retrieve-or-stop decision is recorded at inference time. IRCoT
interleaves retrieval with chain-of-thought steps for multi-step
questions \citep{trivedi2023ircot}. CRAG instead evaluates retrieved
evidence and triggers corrective retrieval actions
\citep{yan2024crag}, while DRAGIN explicitly separates decisions about
when to retrieve from how to formulate the next query
\citep{su2024dragin}. Our sufficiency judge combines the
retrieve-or-stop and follow-up-query decisions. Unlike these systems,
our implementation keeps a fixed generation policy, deterministic
query validation, and explicit document and iteration caps.

Evaluation of retrieve/generate systems has increasingly emphasized
coverage, attribution, and rubric-based assessment. The TREC RAG
workbench and rubric-oriented evaluation studies provide context for
our development diagnostics \citep{dietz2024workbench,farzi2024pencils,pradeep2024initialnugget,pradeep2025greatnugget}.
ALCE evaluates citation-bearing long-form answers along correctness,
fluency, and citation-quality dimensions, and shows that apparently
strong generated answers can still lack complete citation support
\citep{gao2023alce}.
RAGAS also uses automated judgments \citep{es2024ragas}, while work on
faithfulness evaluation examines their interpretation
\citep{jing2025scale}. Separately, citation \emph{correctness} (whether
a cited passage supports a claim) and citation \emph{faithfulness}
(whether the citation reflects genuine grounding rather than
post-hoc rationalization) have been shown to diverge substantially,
with a large share of formally correct citations found unfaithful
under adversarial testing \citep{wallat2025correctness}; our
sentence-level verify/revise stage targets citation correctness. We do
not separately measure faithfulness as defined in that study. The
reported internal diagnostics are development observations and are not
equated with organizer-defined test measures.
